% ============================================================
% Appendix
% ============================================================



\section{Full Proofs}
\label{app:proofs}

This appendix collects the technical material that supports the main
paper's central claims but would be too detailed for the main
narrative. We begin with proof sketches for the sampler properties,
then report the composability analysis, hyperparameter sensitivity, and
per-conversation breakdowns referenced in the main text.

\subsection{Proof of Lemma~\ref{lem:dist-preserve}}
The key idea is the standard binning argument underlying
distribution-preserving keyed selection. Integer-quantize
$\probdist$ into $N$ buckets, apply a cyclic shift determined by a
uniform PRF output, and return the candidate whose interval contains
the shifted point. Marginalizing over the uniform shift recovers the
original interval mass for every candidate, so the keyed pick has the
same marginal distribution as $\probdist$.

\subsection{Proof of Lemma~\ref{lem:cascade}}
Distinct internal decisions use distinct context strings and therefore
distinct nonces $\nonceT = \mathrm{PRF}(K, \ctx)$. Under the standard
PRF assumption, these nonces are computationally indistinguishable from
independent random draws to an adversary without the key. The keyed
picks across a cascade of LLM calls therefore behave independently for
verification purposes, and the expected embedding capacity adds across
valid decision points.

\subsection{Proof of Lemma~\ref{lem:backend-invariant}}
This is an immediate corollary of
Lemma~\ref{lem:dist-preserve}. The marginal expression depends only on
the candidate distribution $\probdist$ and the secret key $K$.
Backend-specific implementation details affect how the candidate set is
enumerated, but once $(\candset, \probdist, \ctx)$ is fixed, the
sampling guarantee is backend-agnostic.

\section{Metrics}
\label{app:metrics}

% ============================================================
% Table: Metric definitions (README §9.2)
% ============================================================
\begin{table*}[t]
\centering
\small
\setlength{\tabcolsep}{3.5pt}
\begin{tabular}{p{0.17\linewidth}p{0.29\linewidth}p{0.46\linewidth}}
\toprule
Category & Metric & Definition / Reporting Unit \\
\midrule

Utility
  & \locomo{} F1
  & precision and recall are computed over predicted versus gold answer
    items, and combined as their harmonic mean. \\

  & BLEU-1
  & Unigram precision between the generated answer and the reference
    answer, measuring lexical overlap at the token level. \\

\midrule

Capacity
  & Bits per memory decision
  & Average number of recoverable watermark bits carried by one write,
    update, delete, link, or merge decision. \\

  & Bits per session
  & Total recoverable watermark capacity within one complete conversation
    session. \\

  & Bits per benchmark episode
  & Total recoverable watermark capacity within one benchmark instance. \\

  & Bits per 100 conversation turns
  & Normalized capacity over long-horizon conversations, reported per
    100 turns. \\

  & Average bits per memory write
  & Average number of recoverable bits carried by each committed memory
    write. \\

\midrule

Carrier-level analysis
  & Per-carrier reliability and capacity
  & All watermark reliability and capacity metrics are additionally
    reported for each carrier:
    $\carrier \in
    \{\text{update\_target}, \text{link\_target},
    \text{semantic\_realization}\}$. \\

\midrule

Robustness
  & Commitment verification fail rate
  & Fraction of tampered memory records whose commitment verification
    fails under the anchored commitment root. Reported per attack type
    in \S\ref{sec:rq4}. \\

  & Merkle proof verification fail rate
  & Fraction of tampered memory records whose Merkle inclusion proof fails
    against the anchored Merkle root. Reported per attack type in
    \S\ref{sec:rq4}. \\

\midrule

Memory integrity
  & Snapshot size
  & Total size of the memory snapshot after watermarking and commitment
    construction. \\

  & Carrier distribution
  & Distribution of embedded bits across carriers, reported overall and
    per $\carrier$. \\

  & Evidence-grounded retrieval recall
  & Fraction of QA-time gold evidence records successfully retrieved from
    memory. \\

  & Write failures
  & Number or fraction of attempted memory writes that fail validation,
    storage, or commitment construction. \\

\bottomrule
\end{tabular}
\caption{Metric definitions across the five research questions. Utility
metrics measure answer quality on \locomo{}; watermark metrics measure
decodability, reliability, and security; capacity metrics quantify how many
bits can be carried by memory operations; carrier-level metrics decompose
performance by watermark carrier; tamper-detection metrics measure whether
memory modifications are caught by commitment and Merkle verification; and
memory-integrity metrics measure the operational side effects of watermarking.
For paper-comparability, F1 follows \amem{}'s set-based formula.}
\label{tab:metric-defs}
\end{table*}
Table~\ref{tab:metric-defs} consolidates the definitions of all metrics
used in this paper, grouped by the research question they address: utility
(RQ1), watermark reliability (RQ2), capacity (RQ3), tamper detection (RQ4),
and memory integrity (RQ5). For each metric we list its reporting unit and,
where relevant, the carrier-level decomposition along
$\carrier \in \{\text{update\_target},\ \text{link\_target},\
\text{semantic\_realization}\}$. We refer the reader to this table
throughout Sections~\ref{sec:rq1}--\ref{sec:rq5} for precise definitions.


\section{Composability Experiment}
\label{app:composability}
We evaluate three configurations --- memory-only, content-only, and
both --- under the R1 protocol of \S\ref{sec:rq3}. The purpose is to
check whether memory-level attribution and output-level watermarking
remain operational when deployed simultaneously. We report
(i) memory decode rate, (ii) content watermark detection, and
(iii) joint utility. Table~\ref{tab:composability} summarizes the
comparison.

% ============================================================
% Table: RQ6 composability (appendix)
% ============================================================
\begin{table}[t]
\centering
\small
\setlength{\tabcolsep}{4pt}
\begin{tabular}{lccc}
\toprule
Configuration & Memory Decode & Content Detect & Joint F1 \\
\midrule
memory-only  & \todo{} & N/A     & \todo{} \\
content-only & N/A     & \todo{} & \todo{} \\
both         & \todo{} & \todo{} & \todo{} \\
\bottomrule
\end{tabular}
\caption{\textbf{RQ6 — Composability with content-layer watermark}
on \S\ref{sec:rq3} R1 protocol. Both watermarks coexist without
mutual interference; joint F1 stays within statistical noise of
either standalone.}
\label{tab:composability}
\end{table}


\section{Hyperparameter Sensitivity}
\label{app:hyperparam}
We sweep $K \in \{2, 4, 8\}$ and
$T_{\text{enum}} \in \{0.5, 0.7, 1.0\}$ to quantify the trade-off
between candidate diversity and embedding stability. We report realized
bits per decision together with the corresponding utility deltas, so
that the reader can see how much attribution capacity is obtained per
unit of behavioral perturbation.

% ============================================================
% Table: Hyperparameter sweep (appendix)
% ============================================================
\begin{table}[t]
\centering
\small
\setlength{\tabcolsep}{4pt}
\begin{tabular}{ccccc}
\toprule
$K$ & $T_{\text{enum}}$ & bits/decision & $\Delta$F1 vs no-wm & R3 Recov. \\
\midrule
2 & 0.7 & \todo{} & \todo{} & \todo{} \\
4 & 0.5 & \todo{} & \todo{} & \todo{} \\
4 & 0.7 & \todo{} & \todo{} & \todo{} \\
4 & 1.0 & \todo{} & \todo{} & \todo{} \\
8 & 0.7 & \todo{} & \todo{} & \todo{} \\
\bottomrule
\end{tabular}
\caption{\textbf{Hyperparameter sweep} on candidate count $K$ and
enumeration temperature $T_{\text{enum}}$. $K{=}4, T_{\text{enum}}{=}0.7$
is the default in the main paper.}
\label{tab:hyperparam}
\end{table}


\section{Per-Conversation Breakdown}
\label{app:per-conv}
Finally, we provide the per-conversation breakdown for all
10 \locomo{} conversations. This table is included only in the
appendix because its role is diagnostic rather than conceptual: it
shows that the aggregate trends in the main text are not driven by a
single conversation outlier.

% ============================================================
% Table: Per-conversation breakdown (appendix)
% Auto-generated from individual conv*_metrics_summary.md files.
% ============================================================
\begin{table*}[t]
\centering
\scriptsize
\setlength{\tabcolsep}{3pt}
\begin{tabular}{c c cccc cccc cccc}
\toprule
& & \multicolumn{4}{c}{\amem{}} & \multicolumn{4}{c}{\graphiti{}} & \multicolumn{4}{c}{Watermark Metrics} \\
\cmidrule(lr){3-6}\cmidrule(lr){7-10}\cmidrule(lr){11-14}
Conv & QAs & {F1 wm} & {F1 nwm} & {$\Delta$} & {Ev. R.} & {F1 wm} & {F1 nwm} & {$\Delta$} & {Ev. R.} & {R1} & {R2 $r{=}0.5$} & {R3} & {Wrong K.} \\
\midrule
0 & \todo{} & \todo{} & \todo{} & \todo{} & \todo{} & \todo{} & \todo{} & \todo{} & \todo{} & \todo{} & \todo{} & \todo{} & \todo{} \\
1 & \todo{} & \todo{} & \todo{} & \todo{} & \todo{} & \todo{} & \todo{} & \todo{} & \todo{} & \todo{} & \todo{} & \todo{} & \todo{} \\
2 & \todo{} & \todo{} & \todo{} & \todo{} & \todo{} & \todo{} & \todo{} & \todo{} & \todo{} & \todo{} & \todo{} & \todo{} & \todo{} \\
3 & \todo{} & \todo{} & \todo{} & \todo{} & \todo{} & \todo{} & \todo{} & \todo{} & \todo{} & \todo{} & \todo{} & \todo{} & \todo{} \\
4 & \todo{} & \todo{} & \todo{} & \todo{} & \todo{} & \todo{} & \todo{} & \todo{} & \todo{} & \todo{} & \todo{} & \todo{} & \todo{} \\
5 & \todo{} & \todo{} & \todo{} & \todo{} & \todo{} & \todo{} & \todo{} & \todo{} & \todo{} & \todo{} & \todo{} & \todo{} & \todo{} \\
6 & \todo{} & \todo{} & \todo{} & \todo{} & \todo{} & \todo{} & \todo{} & \todo{} & \todo{} & \todo{} & \todo{} & \todo{} & \todo{} \\
7 & \todo{} & \todo{} & \todo{} & \todo{} & \todo{} & \todo{} & \todo{} & \todo{} & \todo{} & \todo{} & \todo{} & \todo{} & \todo{} \\
8 & \todo{} & \todo{} & \todo{} & \todo{} & \todo{} & \todo{} & \todo{} & \todo{} & \todo{} & \todo{} & \todo{} & \todo{} & \todo{} \\
9 & \todo{} & \todo{} & \todo{} & \todo{} & \todo{} & \todo{} & \todo{} & \todo{} & \todo{} & \todo{} & \todo{} & \todo{} & \todo{} \\
\midrule
\textbf{Mean} & \todo{} & \todo{} & \todo{} & \todo{} & \todo{} & \todo{} & \todo{} & \todo{} & \todo{} & \todo{} & \todo{} & \todo{} & \todo{} \\
\bottomrule
\end{tabular}
\caption{\textbf{Per-conversation breakdown on \locomo{} all 10
samples}. F1 wm / nwm = \memmark{} vs \texttt{no-watermark} F1 (set
formula). Ev. R. = evidence-grounded retrieval recall. Watermark
metrics use the \amem{} backend numbers.}
\label{tab:per-conversation}
\end{table*}

